\section*{C8: Lambda-calcul cu tipuri simple și PCF}

\subsection*{1. Deducerea tipului}

Considerăm
\[
	T=(\lambda z.\lambda u.z)(yx).
\]
Tiparea principală a variabilelor libere are forma
\[
	\Gamma=\{x:\rho,\ y:\rho\to\beta\}.
\]
Pentru o variabilă de tip nou $\gamma$, obținem
\[
\frac{
	\frac{
		\frac{\Gamma\cup\{z:\beta,u:\gamma\}\vdash z:\beta}
		{\Gamma\cup\{z:\beta\}\vdash\lambda u.z:\gamma\to\beta}
	}{
		\Gamma\vdash\lambda z.\lambda u.z:
		\beta\to(\gamma\to\beta)
	}
	\qquad
	\frac{
		\Gamma\vdash y:\rho\to\beta
		\qquad
		\Gamma\vdash x:\rho
	}{
		\Gamma\vdash yx:\beta
	}
}{
	\Gamma\vdash(\lambda z.\lambda u.z)(yx):\gamma\to\beta
}.
\]
În exemplul à la Church din curs se ia
$\rho=\alpha\to\alpha$, deci
$x:\alpha\to\alpha$ și
$y:(\alpha\to\alpha)\to\beta$.

\subsection*{2. Relația dintre prezentările Church și Curry}

În prezentarea Church, variabilele sunt alese din mulțimi disjuncte,
câte una pentru fiecare tip. În exemplu avem intrinsec
\[
	x:\rho,\qquad y:\rho\to\beta,\qquad
	z:\beta,\qquad u:\gamma,
\]
iar termenul este bine format prin construcție. În prezentarea Curry,
termenul nu conține tipuri; aceleași informații sunt furnizate de
context și de derivarea $\Gamma\vdash T:\gamma\to\beta$. Uitarea
tipurilor intrinseci trece de la Church la Curry, iar o derivare Curry
determină o atribuire Church compatibilă.

\subsection*{3. PCF: adunarea}

Notăm
\[
	\mathsf{add}:=
	\mathsf{fix}(f,\lambda x.\lambda y.
	\mathsf{ifz}(y,x,w,\mathsf{s}((fx)w))).
\]
Regulile de tipare dau
\[
	\vdash\mathsf{add}:\mathbb N\to(\mathbb N\to\mathbb N).
\]
Într-adevăr, în contextul
$f:\mathbb N\to\mathbb N\to\mathbb N$ și $x,y:\mathbb N$, ambele ramuri
ale lui $\mathsf{ifz}$ au tipul $\mathbb N$.

Pentru valori $M,N:\mathbb N$, regulile eager dau
\[
	\mathsf{add}\ M\ \mathsf z\to^* M
\]
și
\[
	\mathsf{add}\ M\ \mathsf{s}(N)
	\to^*\mathsf{s}(\mathsf{add}\ M\ N).
\]
Acestea sunt exact ecuațiile recursive
$m+0=m$ și $m+(n+1)=(m+n)+1$.

\subsection*{4. PCF: scăderea cu punct}

Notăm
\[
	\mathsf{monus}:=
	\mathsf{fix}(f,\lambda x.\lambda y.
	\mathsf{ifz}(y,x,w,\mathsf{ifz}(x,\mathsf z,u,(fu)w))).
\]
Din aceleași reguli,
\[
	\vdash\mathsf{monus}:\mathbb N\to(\mathbb N\to\mathbb N).
\]
Pentru valori $M,N:\mathbb N$:
\[
	\begin{aligned}
		\mathsf{monus}\ M\ \mathsf z
		&\to^* M,\\
		\mathsf{monus}\ \mathsf z\ \mathsf{s}(N)
		&\to^*\mathsf z,\\
		\mathsf{monus}\ \mathsf{s}(M)\ \mathsf{s}(N)
		&\to^*\mathsf{monus}\ M\ N.
	\end{aligned}
\]
Prin inducție după al doilea argument, termenul reprezintă funcția
\[
	m\mathbin{\dot{-}}n=
	\begin{cases}
		m-n,&m\geq n,\\
		0,&m<n.
	\end{cases}
\]
